% ==========================================
% GENERAL INTRODUCTION
% ==========================================
\chapter*{General Introduction}
\markboth{General Introduction}{General Introduction}
\addcontentsline{toc}{chapter}{General Introduction}

With the rapid expansion of network infrastructures, cybersecurity has become central to an organisation's resilience in the digital age. Firewalls are the first line of defence and must be properly managed and configured to safeguard sensitive data from the ever-evolving threat landscape. Traditional methods of firewall administration, which rely on manual command-line input and complex graphical user interfaces, cannot meet the speed and complexity demanded by modern networks.

The most significant challenge to network administration today is the "semantic gap" between high-level security goals and low-level technical implementation. The security administrator's intent must be translated into vendor-specific syntax (i.e., configuration rules) manually, and this is not only a lengthy process but also very error-prone. A small misconfiguration of a security policy can result in a serious vulnerability and/or a disruption of service.

This final-year project, carried out at \textbf{Kleos Tunisia}, addresses these inefficiencies at the intersection of Artificial Intelligence and Network Security. The main goal is to develop and implement a hybrid AI-assisted administration agent for \textbf{Fortinet FortiGate} firewalls. We propose a system that leverages Large Language Models (LLMs) and Retrieval-Augmented Generation (RAG) to enable administrators to interact with infrastructure through natural language.

The innovation of the project lies in its multi-layered approach:
\begin{itemize}
    \item \textbf{Natural Language Processing:} Understanding administrative intent in English and French with the Mistral model.
    \item \textbf{Grounding in context:} Implement a RAG mechanism based on ChromaDB to ensure that all generated configurations strictly follow the official Fortinet technical documentation.
    \item \textbf{Security and Verification:} A rigorous validation pipeline and human-in-the-loop mechanism will be in place to ensure no configuration is deployed without validation and structured audit logging.
\end{itemize}

In this work, we show how generative AI can be transformed from a general-purpose tool into a specialized, trustworthy, and safe network administration assistant, thereby reducing operational overhead and improving the enterprise's overall security posture.
